\section{Introduction}
\label{sec:intro}

I/O data passes through the kernel on its way between an application and a device.
A proxy may forward that data unchanged, yet ordinary socket I/O copies it into
application memory and back, using CPU time and memory bandwidth~\cite{cai21,libra}.
Zero-copy sends can avoid the second copy by borrowing application pages. The kernel
uses pinned pages and tells the application when they may be reused~\cite{msgzerocopy,sendzcdoc}.
The loan can outlast device completion because TCP may keep the bytes for retransmission.

Lending receive buffers to an application raises a protection problem. An application
can change any byte in a writable mapping, including packet headers that arrive in
the buffers. Some zero-copy receive interfaces use hardware to place headers in
separate memory and steer flows to application-specific queues~\cite{zcrx,maio}.
We show how software can isolate kernel control state while letting a trusted
application edit I/O data on commodity network devices.

\sys follows the principle \emph{lend, don't borrow}. The kernel owns the pages and
gives the application a handle naming the data. The application decides which bytes to
keep and where to send them. Linux continues to process TCP and UDP, and an object can
outlive the socket or device operation that produced it. The application can release
its handle while outstanding I/O holds references to the pages. The kernel reclaims
the pages when those references are gone. The application still receives I/O results,
and the kernel handles reclamation without a separate buffer-reuse notification.

\paragraph{Named mode.} Applications use handles to read selected bytes, retain objects
and direct later I/O. The kernel checks each operation and leaves the underlying pages
unmapped. Socket capture can use ordinary receive buffers and existing drivers.
Authority to access a socket or block device is separate from authority to use an object.

\paragraph{Mapped mode.} Trusted applications can read and edit buffers in a dedicated
I/O area. Kernel control structures and the header copies used for parsing stay in
private memory. The network driver supplies receive buffers from that area. Both modes
work without hardware header splitting or per-application flow steering. A mapper must be
trusted with the entire area, which can include all traffic on a device. It must keep
bytes stable while outstanding I/O uses them.

TCP, UDP and block I/O illustrate how the same two access modes apply to different
endpoints. A TCP relay can retain a range after a connection closes. A UDP relay handles
one datagram at a time. A storage server can read an aligned range into an object
and send it in several replies. The same model also supports writes to compatible
devices, so a network service can use it in both directions. In each case, the
choice between named and mapped access determines how the application touches the
bytes. The endpoint determines the operation's boundaries, permissions and completion.

IO-Lite and PASTE establish the value of sharing retained I/O data, and MAIO establishes
segregated writable I/O memory~\cite{iolite,paste,maio}. \sys combines checked operations
on unmapped objects with writable objects backed by an area segregated in software.
Both forms share the same retention and reuse model. We describe that model and its
Linux realization (\S\ref{sec:lending}--\ref{sec:impl}), then show how applications can
use it across stream, datagram and block I/O (\S\ref{sec:apps}).
